\begin{textsample}{NEG Dim 9 – global\_north\_2025\_04 – Score -6.00 – global\_north\_2025\_04/123856663.txt}  \label{ex:f9_neg_035}
Follow us In Jerusalem, Christians celebrate Easter at gunpoint Armed police gathered at Al-Aqsa Mosque to escort thousands of Israeli settlers to its courtyard Christians carrying crosses walk through Jerusalem ' s Old City on their way to the Church of the Holy Sepulchre on 18 April 2025 ( Ahmad Gharabli/AFP ) There were barely any Palestinians to be seen in Jerusalem ' s Old City during the Jewish holiday of Passover.
The usual crowds of worshippers that line the main streets leading to Al-Aqsa Mosque were replaced by police officers, armed settlers and metal gates.
Every shop was forced to close.
A Palestinian woman, one of the few I met in the old city, told me :"They Israelis spit on us ; they target those of us who wear scarves."The situation had reached such lows, she said, that she now avoided going to worship at the revered site during any Jewish holiday.
Armed police had amassed at Al-Aqsa to escort thousands of Israeli settlers to the mosque @ @ @ @ @ @ @ @ @ @ other far-right politicians were among them.
New \textbf{MEE} \textbf{newsletter} : Jerusalem Dispatch Sign up to get the latest insights and analysis on Israel-Palestine, alongside Turkey \textbf{Unpacked} and other \textbf{MEE} \textbf{newsletters} For these politicians, the record number of Jewish worshippers at a spot they call the Temple Mount was a major cause for celebration."Historic records on the Temple Mount : A 37 \% surge in the number of pilgrims during Passover, with 6,315 pilgrims entering the Temple Mount for Passover prayers.
An all-time record number,"the Temple Mount administration tweeted.
It was a shocking scene to witness, one of Islam ' s holiest sites filled with sounds usually found in a synagogue.
Anger is clearly bubbling away among Christian and Muslim worshippers who are prevented from practicing their religions This may have been a cause for celebration for the Israeli far right, but it was another deep wound inflicted on Muslims who are repeatedly barred from accessing Al-Aqsa Mosque to worship.
The Islamic Waqf, the Jordanian-controlled organisation responsible for @ @ @ @ @ @ @ @ @ @ irrelevant.
An official who works for the body said they had"zero control over anything".
The official said the body had to wait more than a month to get permission to allow even children ' s toys into Al-Aqsa library.
It is clear that even the writ of King Abdullah II of Jordan, a man who holds the title of official custodian of Al-Aqsa Mosque, has no real sway over what happens at the site.
Carpets donated by the king are still awaiting permission to be brought in, the official said.
But it is crucial to note that the Old City ' s settler overlords do not spare Christians any more than they do Muslims.
Religious event turned into a war zone On the other side of the Old City, Israeli police assaulted Palestinian Christians attempting to access the Church of the Holy Sepulchre.
One member of the border police, an anti-terrorist force, was seen aiming his gun at a member of the Greek Orthodox Scouts.
The scenes of @ @ @ @ @ @ @ @ @ @ turned a day of religious celebration into a war zone.
Even Apostolic Nuncio Adolfo Tito Yllana was not spared the humiliation of being held at a barrier and refused permission to attend Easter ceremonies at the Church of the Holy Sepulchre.
One Christian worshipper who managed to run the gauntlet of the police and get through to the church posted a video showing an empty square - not one that Israeli authorities claimed was dangerously overcrowded.
Huge surge of Jewish worshippers at Al-Aqsa Mosque as Muslims locked out Since the coronavirus pandemic, Israel has imposed strict restrictions on the number of Christians who can enter the church and the square during holidays.
They claim the restrictions are for the public ' s own safety and good.
Yanal Jabarin, a Palestinian journalist, says that the restrictions are getting worse and worse every year.
Anger is clearly bubbling away among Christian and Muslim worshippers who are prevented from practising their religions.
Xavier Abu Eid, a political scientist and former adviser to the Palestine Liberation Organisation @ @ @ @ @ @ @ @ @ @ plan to change the status quo of Jerusalem ' s holy sites."And if they succeed, I don't think there will be much future for Palestinian Christians in the holy city,"he said."They are trying to empty the Christian Quarter and Armenian Quarter."' Easter in the shadow of genocide ' When Christians were allowed to celebrate Easter this year, the mood was sombre.
In Gaza, like much of the West Bank, traditional Easter festivities were replaced by quiet church services and symbolic gestures of resilience.
Palestinian Christians marked a solemn Easter in northern Gaza, as Israeli forces continued to lay siege to the enclave ' s largest hospital and bombed several targets across the embattled territory.
Reverend Munther Isaac, a prominent Christian pastor from Bethlehem, warned that Christians in both the occupied West Bank and the war-battered Gaza Strip were facing potential extinction amid relentless Israeli assaults."For the second year in a row, we celebrate Easter in the shadow of @ @ @ @ @ @ @ @ @ @ the Anadolu news agency last week."Palestine is still walking the path of sorrow, suffering from Israeli siege and apartheid policies,"he said."The same violence that killed Christ still exists in our land today."The views expressed in this article belong to the author and do not \textbf{necessarily} reflect the editorial policy of Middle East Eye.
Lubna Masarwa is a journalist and Middle East Eye ' s Palestine and Israel bureau chief, based in Jerusalem.
Middle East Eye delivers independent and \textbf{unrivalled} coverage and analysis of the Middle East, North Africa and beyond.
To learn more about \textbf{republishing} this content and the associated fees, please fill out this form.
More about \textbf{MEE} can be found here.

% matched lemmas: mee, necessarily, newsletter, republish, unpack, unrivalled
\end{textsample}
